\subsection{The work behind the checkpoints}

A project owner assembles a dossier describing the proposed change: its need, users, budget,
technical design, supplier commitments, and readiness evidence. A specialist gate examines the
parts relevant to its policy and returns a recorded opinion. The same dossier can therefore
receive different findings from Architecture, Security, and Legal without those reviews being
duplicates. General combines their effective opinions with its own governance checks.
Table~\ref{tab:gate-guide} explains the eight families represented in the experiment.

\begin{table}[!htb]
\centering\small
\caption{What each benchmark gate reviews. Questions summarize the supplied policies; listed
records and checks are examples, not exhaustive definitions of the professions.}
\label{tab:gate-guide}
\begin{tabularx}{\linewidth}{@{}P{2.45cm}LL@{}}
\toprule
Gate & Review question & Examples of material examined\\
\midrule
Procurement & Does the selected supplier satisfy purchasing requirements?
& Offers, mandatory criteria, due diligence, sanctions, and three-year cost\\
Legal & Are the contractual prerequisites and signing authority present?
& Data-processing agreement, liability, exit assistance, and signing authority\\
Compliance & Does the proposal satisfy the applicable compliance checks?
& Required impact assessment, residency, regulatory mapping, and audit trail\\
Security & Does the design provide the protections required for its exposure and data?
& Identity controls, private endpoints, monitoring, and vulnerability findings\\
IT & Does the proposal fit the technology and operating environment?
& Catalog status, existing capability, capacity, support ownership, and change records\\
Architecture & Is the proposed system consistent with the declared design constraints?
& Network overlap, interfaces, data ownership, latency, and reversibility\\
Tech Readiness & Is the evidence of operational readiness sufficient?
& Load, restore and recovery tests, runbook, rollback, and handover\\
General & What governance disposition follows from the project and specialist reviews?
& Effective upstream opinions, budget, strategic alignment, and change plan\\
\bottomrule
\end{tabularx}
\end{table}

The five dispositions express different completed review outcomes. \texttt{GO} permits
proceeding within the reviewed scope. \texttt{GO\_WITH\_RESERVATIONS} retains conditions
alongside that permission. \texttt{REWORK} requests changes to the submitted proposal;
\texttt{SUSPENSION} waits for a prerequisite; \texttt{NO\_GO} rejects proceeding as submitted.
The policy and mandate determine which outcome is admissible. The experimental schedule
collects every review in the route, including after a blocking opinion, to evaluate the
remaining specialist work and final consolidation.

Four responsibilities clarify the unit of substitution: preparing the dossier, examining
it, making an authorized decision effective, and carrying out the resulting actions.
An organization may allocate several of these responsibilities to one person, or distribute
them across project teams, specialists, and decision owners. In the experiment, the agent
performs the specified review and can use authorized simulated actions. The generator supplies
the dossier; the trace records requested remediation and commitments. A deployment's labor
account follows each responsibility wherever it is performed.

For example, a readiness review may identify a draft runbook, require its completion, and
record an authorized conditional decision. The review has then produced its required output;
the runbook still has an owner and an open completion condition. The Falcon trace in
Section~\ref{sec:benchmark} records precisely this separation. It explains how an agent can
replace a review task while another person or system continues the operational work triggered
by that decision. Extending substitution to that operational work requires its own accepted
output and authority.
